\section{Evaluation Results v1.4--v1.8: Quantitative Metrics and Version-by-Version Progress}
\label{sec:evaluation}

\subsection{Evaluation Philosophy: From Output Validation to Behavioral Verification}

UMAF's evaluation infrastructure evolved through four phases across versions v1.4--v1.8, growing from 8 unit tests to 379 behavioral tests across 10 files. This growth represents a philosophy shift from \textbf{output validation} (checking that pipelines produce files) to \textbf{behavioral verification} (testing that graph nodes, parse\_result logic, flow routing, fallback methods, and resume state reconstruction all behave correctly) \cite{evaluatingcompound}.

\subsection{Test Infrastructure Architecture}

The 10-file modular structure:
\begin{itemize}
    \item \texttt{conftest.py}: Shared fixtures (MockLLM, tmp\_working\_dir, mock\_agent\_result)
    \item \texttt{test\_smoke.py}: 42 tests for pipeline instantiation and graph topology
    \item \texttt{test\_pipeline.py}: BasePipeline utilities (topological\_levels, run\_workers\_with\_deps, status\_router)
    \item \texttt{test\_coder.py}: 27 tests (decomposition, review loop, router logic, coder\_files injection)
    \item \texttt{test\_research.py}: 62 tests (decomposition, worker retry, version bump, reviewer scoring, writer LaTeX)
    \item \texttt{test\_coderpp.py}: 58 tests (decomposition, worker deps, reviewer loop, observer, organizer)
    \item \texttt{test\_topology.py}: 14 tests (analyzer, designer, evaluator, writer roles; fallback chains)
    \item \texttt{test\_skill.py}: 15 tests (scanner, 4 detectors, aggregator, writer; artifact classification)
    \item \texttt{test\_feature.py}: Tests for scanner, planner, coder, reviewer roles
    \item \texttt{test\_self\_evolution.py}: 49 tests (analyzer, planner, coder, reviewer, writer; git diff detection; test suite integration)
\end{itemize}

\subsection{Behavioral Test Design}

v1.8 behavioral tests verify causal chains rather than structural properties:
\begin{itemize}
    \item \textbf{Graph node behavior}: Each pipeline node correctly transforms state with controlled inputs
    \item \textbf{parse\_result logic}: Tested with three states (valid LLM output, disk fallback, fallback code path)
    \item \textbf{Flow routing}: All defined statuses produce correct transitions; unknown statuses route to END
    \item \textbf{Fallback methods}: Independent LLM-free testing; verify structurally valid output
    \item \textbf{Resume state reconstruction}: Checkpoint states at each pipeline stage
    \item \textbf{Cross-key dependency resolution}: Both int IDs and string module names resolve correctly
\end{itemize}

\subsection{Version-by-Version Quantitative Trajectory}

\begin{table}[H]
\centering
\small
\caption{Test Growth Across Versions}
\label{tab:testgrowth}
\begin{tabular}{@{}lccccc@{}}
\toprule
\textbf{Version} & \textbf{Tests} & \textbf{Files} & \textbf{Pipelines} & \textbf{Roles} & \textbf{Key Event} \\
\midrule
v1.3   & 8   & 1  & 3 & 8  & Baseline (unit tests) \\
v1.4.1 & 23  & 1  & 3 & 10 & +15 agent loop smoke tests \\
v1.5   & 42  & 1  & 5 & 20 & +19 pipeline smoke tests \\
v1.6   & 99  & 5  & 6 & 25 & Structural split into 5 files \\
v1.6.1 & 131 & 5  & 6 & 25 & +32 DI fix tests \\
v1.7   & 99  & 5  & 6 & 25 & Revalidated after dedup \\
v1.8   & 379 & 10 & 7 & 32 & +280 tests (+175 behavioral + 49 SelfEvolution) \\
\bottomrule
\end{tabular}
\end{table}

\subsection{Production Metrics}

\textbf{Worker Success Rate} (ResearchPipeline): v1.3: 4/6 (66.7\%) $\rightarrow$ v1.4: 7/7 (100\%). The improvement was driven by two mechanisms: the 600s timeout giving workers adequate time, and the \texttt{parse\_result()} honesty fix.

\textbf{Review Scores} (ResearchPipeline, /50): v1.3: range 43--31, mean \textasciitilde37.5 $\rightarrow$ v1.4: range 48--38, mean \textasciitilde43.3.

\textbf{LaTeX Output Size}: v1.3: 41KB $\rightarrow$ v1.4: 60KB with sectioned output and \texttt{\textbackslash input\{\}} commands.

\textbf{Pipeline Execution Time} (ResearchPipeline): v1.3: \textasciitilde360s $\rightarrow$ v1.4: 443s with full 7-worker coverage.

\textbf{Codebase Metrics}: v1.4: \textasciitilde3,500 lines $\rightarrow$ v1.8: \textasciitilde8,500 lines. Test density: \textasciitilde1 test per 22 lines of framework code.

\subsection{Dual-Layer Evaluation Architecture}

UMAF's novel contribution \cite{multiagentbench,collabovercooked} is a dual-layer design where the \textbf{online layer} (reviewer agents using REVIEW\_PASSED/REVIEW\_FAILED token scanning) gates pipeline progress during execution, and the \textbf{offline layer} (pytest suite) verifies the correctness of the online layer. This creates a recursive verification structure \cite{reflexion,selfrefine}: the test suite verifies that \texttt{scan\_review\_verdict()} correctly parses verdict tokens, and \texttt{scan\_review\_verdict()} gates whether pipelines produce valid output.

\subsection{v1.4.1 Agent Loop Reliability Foundations}

The v1.4.1 bug fixes deserve detailed analysis as the reliability foundation for all subsequent improvements:

\raggedright
\textbf{Tool Execution Before TASK\_COMPLETE}: The original loop checked TASK\_COMPLETE before executing tools, silently dropping tool calls that accompanied completion declarations. The fix reorders: execute tools first, then check TASK\_COMPLETE only when no pending tool calls remain.

\textbf{Error Spiral Threshold 3$\rightarrow$2}: Empirical analysis of 127 agent logs found that of 31 agents reaching 3 consecutive errors, only 6.5\% self-corrected. Of 47 reaching 2 consecutive errors, 19.1\% self-corrected at iteration 3. Threshold 2 saves \textasciitilde1 iteration per spiral while catching 93.5\% of doomed agents.

\textbf{Post-Loop File Verification}: After loop exit, checks whether agent produced any output file. If no file exists but write\_file calls were made, replays the last attempted write. Catches cases where JSON escaping errors cause silent write failures.
